% Cover letter and point-by-point reviewer responses.
% Reviewer text is verbatim from genetics_review.txt; \reply{...} blocks
% are scaffolded as TODO placeholders to be filled in during the revision.
% Wherever a change is made in the manuscript in response to a point,
% mark the manuscript line with \revpoint{R}{P} and cite it here with
% \revref (inside the matching {point}) or \revreffull{R}{P}.
%
% Reviewer labels used for \revpoint: E (editors), 1, 2.

\begin{minipage}[b]{2.5in}
  Resubmission Cover Letter \\
  {\it GENETICS} \\
  Manuscript GENETICS-2026-309883
\end{minipage}
\hfill
\begin{minipage}[b]{3.0in}
  \raggedleft
  Angel G.\ Rivera-Col\'on (on behalf of all authors) \\
  \today
\end{minipage}

\vskip 2em

\noindent
{\bf Dear Dr.\ Larracuente,}

\vskip 1em

TODO

\vskip 2em

\noindent\hspace{4em}
\begin{minipage}{4.5in}
\noindent
{\bf Sincerely,}

\vskip 2em

{\bf
  Angel G.\ Rivera-Col\'on (on behalf of all authors)
}\\
\end{minipage}

\vskip 4em

\pagebreak

%%%%%%%%%%%%%%%%%%%%%%%%%%%%%%%%%%%%%%%%%%%%%%%%%%%%%%%%%%%%%%%%%%%%%%%%%
% Editors' comments. Set up by hand rather than with \reviewersection so
% the heading doesn't read "Reviewer E"; points are tagged \revpoint{E}{P}.
\renewcommand{\thereviewer}{E}
\setcounter{point}{0}
\section*{Editors:}

\begin{itquote}
Two experts in the field have reviewed your manuscript, and the editors
have read it as well. As you can see below, both reviewers found your
manuscript well written and your system interesting, but they vary in
their view of the impact. Each reviewer has specific concerns that you
can read at the end of this e-mail. Your work clearly presents a valuable
empirical resource, however the current manuscript does not adequately
place the results in the context of other work in the field. While your
manuscript is not currently acceptable for publication in GENETICS, we
would welcome a substantially revised manuscript. Both reviewers and the
editors have comments and concerns to be addressed in a revised
manuscript.

If you do submit a revised manuscript, the revisions will need to clearly
place this work in a broader context. Please carefully consider the
comments from reviewer 2 as your revise. To be further considered by
GENETICS, the revised manuscript should emphasize how this study
contributes to our understanding of how population size interacts with
population genetic mechanisms. The revised manuscript should include
additional analyses and a major revision to the text.
\end{itquote}

\reply{TODO}

% Editor point #1
\begin{point}{} % dN/dS vs MK tension
There is no need to discuss a tension between dn/ds and the MK tests, as
you explain in the discussion that genes can experience both purifying
and positive selection.
\end{point}

\reply{TODO}

% Editor point #2
\begin{point}{} % purifying vs background selection language
Please carefully revise your language about purifying selection and
background selection, as noted by reviewers. Only use background
selection if you are are referring to the effects of purifying selection
on linked variation.
\end{point}

% TODO: changed them add text
\reply{TODO}

% Editor point #3
\begin{point}{} % TE adaptation too speculative; TE age / size distributions
The discuss of TE adaptation is too speculative based on the data that
you present. A deeper analysis of TE age and size distributions may help
you distinguish between some of your alternative hypotheses.
\end{point}

\reply{TODO}

% Editor point #4
\begin{point}{} % focus on 4-fold pi rather than genomic pi
Don't focus so much on genomic pi. Many factors can influence this,
especially the proportion of repetitive and/or non-functional DNA in a
genome. It is much better for comparative work to focus on 4-fold
degenerate pi, as these sites are easily alignable and under weaker
selection in most lineages.
\end{point}

\reply{TODO}

% Editor point #5
\begin{point}{} % repeat-masked reference for SNP calling?
Please clarify if you used a repeat-masked reference for SNP calling, as
this may affect the level of polymorphism in non-coding regions
especially given the high repeat content of your genome.
\end{point}

\reply{TODO}

% Editor point #6
\begin{point}{} % quantitative context: 4-fold pi, Ne, MK
As noted by reviewer 2, it is important that you make an effort to
quantitatively place your findings in the context of the broader
literature addressing the relationship between 4 fold degenerate pi, Ne,
and MK tests.
\end{point}

\reply{TODO}

% Editor point #7
\begin{point}{} % correlation between dN/dS and synonymous pi
It might be interesting to look at the correlation between dN/dS and
synonymous pi. If directional selection dominates, then the correlation
might be negative.
\end{point}

\reply{TODO}

% Editor point #8
\begin{point}{} % line 47: Hill-Robertson effects
Line 47 - better stated as Hill Robertson effects
\end{point}

\reply{TODO}

% Editor point #9
\begin{point}{} % shared ortholog MK results
It may be interesting to look for shared ortholog MK results, if there
are such comparative resources.
\end{point}

\reply{TODO}

% Editor point #10
\begin{point}{} % Fig 5d: low pi regions / centromeres, physical scale
In Fig 5d, regions of low pi are indicative of centromeres exhibiting
reduced crossing-over. What is the physical scale of the phenomenon? It's
hard to tell from the figure. Are there properties of the variation or
divergence indicative of weaker selection in such regions?
\end{point}

\reply{TODO}

% Editor point #11
\begin{point}{} % lines 908-909: source of mutation rate
Lines 908-909, Please clarify where this mutation rate comes from in the
text. It appears to reference an average of 4 arthropod species rather
than one estimated based on S. balanoides data, as described in the Nunez
PNAS paper.
\end{point}

\reply{TODO}

% Editor point #12
\begin{point}{} % B. crenatus window pi from assembly heterozygosity
Regarding the B. crenatus assembly. Could one estimate a window pi from
heterozygosity in the individual used for the assembly and then
investigate whether the orthologous window pi is correlated in the two
barnacle species?
\end{point}

\reply{TODO}

% Editor point #13
\begin{point}{} % line 954: 'moderate evidence' of non-compositional ENC bias
Line 954, Be explicit about the evidence you are alluding to here
regarding the `moderate evidence' of non-compositional bias in the ENC
reduction.
\end{point}

\reply{TODO}


%%%%%%%%%%%%%%%%%%%%%%%%%%%%%%%%%%%%%%%%%%%%%%%%%%%%%%%%%%%%%%%%%%%%%%%%%
\reviewersection{1}

\begin{itquote}
This is a very strong manuscript that tackles important questions in
evolutionary biology. I agree with the authors that species with
exceptionally large Ne offer unique opportunities to investigate the
efficacy of natural selection and other fundamental evolutionary
processes, making barnacles ideal systems to tackle these questions.

The authors provide a chromosome-level genome assembly and pair this
resource with comprehensive molecular evolutionary analyses that paint a
compelling picture of the strong efficacy of selection in the species.
Overall, I am excited about this paper. I only have a few comments for
the authors to consider:
\end{itquote}

\reply{TODO}

% Reviewer 1 point #1
\begin{point}{Purifying vs.\ Background selection}
I found the evidence for efficient purifying selection in B. glandula
convincing. Yet, I got the impression that the terms ``purifying
selection'' and ``background selection'' were being used somewhat
interchangeably. To me, the manuscript makes a very strong case for the
former, whereas the latter feels harder to establish without additional
information on recombination and patterns of linkage. Since the authors
note that characterizing recombination will be focus of future work, I
encourage caution when invoking background selection in the absence of
those data.
\end{point}

% TODO: changed them add text
\reply{TODO}

% Reviewer 1 point #2
\begin{point}{Adaptive TEs}
The possibility that transposable elements contribute to adaptation in
this species is certainly intriguing. However, I would like to see a
somewhat broader discussion of non-adaptive explanations. The manuscript
paints a very consistent picture of highly efficient selection, and given
those patterns, one has to wonder why (or how) is such TE-rich genome
maintained... could there also be a role for unique dynamics of
insertion-deletion balance, or maybe heterogeneity in the recombination
landscape, etc?
\end{point}

\reply{TODO}

% Reviewer 1 point #3
\begin{point}{Codon bias}
I think the ENC/GC composition analyses are very interesting. Given the
high GC-rich nature of the genome, could GC-biased gene conversion (gBGC)
also have contributed to the observed patterns of codon usage??
\end{point}

\reply{TODO}

% Reviewer 1 point #4
\begin{point}{The balancing selection outlier}
One result that caught my attention was the candidate locus exhibiting a
strong excess of polymorphism relative to divergence. I realize this
observation is not a central focus of the manuscript, and I am not
suggesting additional analyses. However, I was curious whether the
authors have any additional thoughts about this locus. The observation
that this homolog has been implicated in osmotic stress responses makes
the result particularly intriguing, and I wonder whether the authors view
this pattern as potentially consistent with ecological processes capable
of maintaining polymorphism (e.g., intertidal zonation like that of Mpi).
I ask this simply out of curiosity.
\end{point}

\reply{TODO}


%%%%%%%%%%%%%%%%%%%%%%%%%%%%%%%%%%%%%%%%%%%%%%%%%%%%%%%%%%%%%%%%%%%%%%%%%
\reviewersection{2}

\begin{itquote}
This manuscript presents a new chromosome-level genome assembly and a
population genomic analysis of the Pacific acorn barnacle Balanus
glandula, a species characterized by very large census population sizes
and exceptionally high genetic diversity. The authors report genome-wide
nucleotide diversity exceeding 5\%, investigate the distribution of
diversity across functional categories, patterns of codon usage,
demographic history, and adaptive protein evolution. They find reduced
diversity in functional elements, as expected under purifying selection,
and infer a high contribution of positive selection to protein
divergence. Overall, the analyses are carefully conducted and the
manuscript is clearly written. B. glandula is certainly an interesting
addition to the relatively small set of species for which such high
levels of nucleotide diversity have been characterized at the
whole-genome level.

I nevertheless have two major concerns, one concerning the scope of the
study and the other its positioning with respect to the existing
literature.
\end{itquote}

\reply{TODO}


% --- Major concerns ---

% Reviewer 2 point #1
\begin{point}{Scope}
Although I found the analyses technically sound, I am not convinced that
the results provide sufficient conceptual advance.

The main novelty is essentially the detailed genomic characterization of
a single, exceptionally polymorphic species. This is valuable, but highly
polymorphic species have already been included in broad comparative
studies of genetic diversity and molecular evolution, including species
with nucleotide diversity in the same general range. The manuscript
itself cites some of this work, including Cutter's work on highly
polymorphic Caenorhabditis, but the comparative context is narrower than
it should be. Other classical examples include Ciona savignyi,
Caenorhabditis remanei and C. brenneri, and Schizophyllum commune. Large
comparative datasets such as Romiguier et al. (2014) and subsequent
studies by Galtier (on Codon Usage Bias, cited in the ms, or on adaptive
protein evolution, not cited in the ms
\url{https://doi.org/10.1371/journal.pgen.1005774}) have also included
highly polymorphic non-model species (with pi$>$5\%) and investigated the
relationships between genetic diversity, effective population size,
efficacy of purifying selection and adaptive protein evolution.

Consequently, most of the major conclusions here are consistent with
expectations and previous empirical work: very high diversity in a
species with a very large population size, stronger signatures of
purifying selection in functional sequences, and efficient adaptive
protein evolution, although Lewontin's paradox remains. I see
considerable value in documenting these patterns in B. glandula,
particularly with a high-quality genome, but less of a general conceptual
advance than suggested by the manuscript.

I therefore think that this is a good population genomic study, but one
whose contribution is primarily the genomic characterization of an
interesting non-model species.
\end{point}

\reply{TODO}

% Reviewer 2 point #2
\begin{point}{The marine population genetics literature is largely overlooked}
Another concern is the literature review and, consequently, the framing
of the study. I was quite surprised by how little the manuscript engages
with the extensive literature on genetic diversity and population
genetics of marine invertebrates.

Marine invertebrates have been central to discussions of extreme genetic
diversity since the allozyme era. High polymorphism in highly fecund
marine organisms, its relationship with census and effective population
size, and the unusual demographic processes affecting such species have
been studied for decades. In this context, I find it particularly
surprising that the work of Hedgecock and others is not discussed. The
concepts of sweepstakes reproductive success, very large variance in
reproductive success, low Ne/N ratios and, more recently, multiple-merger
coalescents and chaotic genetic patchiness are directly relevant to the
biological interpretation of the patterns reported here.

This omission matters because several observations presented as
characteristic consequences of the exceptional population size of B.
glandula have precedents in marine population genetics. For example, high
nucleotide polymorphism and its consequences for protein variation have
been investigated in marine bivalves for many years. Sauvage et al.
(2007) explicitly studied both nucleotide polymorphism and codon usage in
the Pacific oyster, reporting levels of diversity among the highest known
in animals at the time. Harrang et al. (2013) investigated the unusually
high load of non-neutral amino-acid polymorphisms in the European flat
oyster and explicitly discussed its relationship with high fecundity and
sweepstakes reproduction. More generally, marine genomes are now known to
contain extensive structural variation. Massive gene presence-absence
variation has been described in mussels (e.g. Gerdol et al. 2020), and
abundant CNVs and other structural variants have been documented in
oysters (Modak et al. 2021, Lui et al. 2025) and other marine
invertebrates (e.g. tunicates: Satou et al. 2021, lobster: Dorant et al.
2020). Thus, the genomic complexity observed here should also be
discussed in the context of this literature rather than primarily in
comparison with classical model organisms.

Similarly, mutation rate is no longer entirely unexplored in highly
fecund marine invertebrates. Direct germline mutation-rate estimates are
available, for example in sea stars (Popovic et al. 2024), and relevant
estimates also exist in tunicates. These studies are particularly
relevant to the discussion of Lewontin's paradox because they explicitly
address the respective contributions of mutation rate, census population
size and demographic history to high genetic diversity.

I am not suggesting that all these earlier studies addressed exactly the
same questions or had access to chromosome-level whole-genome data.
Clearly they did not. Rather, my concern is that the manuscript currently
creates the impression that highly polymorphic species, and particularly
highly fecund marine invertebrates, represent a largely unexplored
empirical domain. This is not the case. A more thorough treatment of this
literature would both provide a fairer historical perspective and make it
clearer what the present study genuinely adds.
\end{point}

\reply{TODO}


% --- Additional comments ---

% Reviewer 2 point #3
\begin{point}{} % codon usage: preferred vs unpreferred codons (Fop), gBGC
I am not fully convinced by the codon-usage analysis. ENC is useful as a
general measure of codon bias, but it does not directly identify which
synonymous codons are selectively preferred. In the context of this
manuscript, an analysis based on preferred versus unpreferred codons, for
example using Fop or a related measure after identifying preferred
codons, would be considerably more informative. This is particularly
important because GC3 correction alone does not fully resolve the
distinction between selection on codon usage and GC-biased gene
conversion or other GC-associated processes. Knowing whether preferred
codons are GC- or AT-ending would help discriminate among these
explanations. The literature on codon usage in highly polymorphic
species, including comparative studies by Galtier and colleagues and the
earlier work in Pacific oyster by Sauvage et al. (2007), would provide
useful context here.
\end{point}

\reply{TODO}

% Reviewer 2 point #4
\begin{point}{} % interpreting high alpha; Galtier 2016
Relatedly, I would be cautious about interpreting high alpha from
McDonald-Kreitman analyses too directly as evidence that the
exceptionally large population size of B. glandula produces unusually
strong adaptive evolution. Comparative analyses have already shown that
the relationship between effective population size and alpha is not
straightforward. In particular, Galtier (2016) found that much of the
relationship between diversity and alpha across animals could be
explained by more efficient purifying selection and a reduced
contribution of slightly deleterious substitutions, rather than by a
higher rate of adaptive substitution per se. This distinction seems
particularly relevant here.
\end{point}

\reply{TODO}

% Reviewer 2 point #5
\begin{point}{} % sweepstakes reproduction / multiple-merger coalescents
The demographic interpretation could also engage more explicitly with
reproductive biology. For a highly fecund broadcast-spawning marine
species, a Wright-Fisher model based on conventional Ne should at least
be discussed alongside sweepstakes reproduction and multiple-merger
genealogies. These processes can substantially alter the relationship
between nucleotide diversity, census size and the demographic parameters
inferred under a standard Kingman coalescent.
\end{point}

\reply{TODO}

% Reviewer 2 point #6
\begin{point}{} % connect to existing B. glandula popgen / latitudinal cline
Finally, given that B. glandula itself has a substantial population
genetic and phylogeographic literature, including work on its strong
latitudinal genetic cline, I think the manuscript could do more to
connect the new genome-wide results to what is already known about
population structure and demographic history in this particular species.
\end{point}

\reply{TODO}

% Reviewer 2 point #7
\begin{point}{Overall assessment}
This is a carefully conducted study based on a valuable new genomic
resource, and I have no major concerns about the overall quality of the
analyses. My reservations concern primarily the level of novelty claimed
and the incomplete treatment of previous work. The study provides an
interesting whole-genome characterization of an extremely polymorphic
marine species, but many of its main biological conclusions have
precedents in comparative population genomics and, especially, in the
long-standing literature on highly fecund marine invertebrates. A
substantial revision of the introduction and discussion would therefore
be necessary to place the results in their appropriate context and to
more clearly identify what is genuinely new.
\end{point}

\reply{TODO}
